\documentclass[10pt, aspectratio=169]{beamer}
\input{../../_en_preambulo_beamer}
% Comandos y macros estadísticas compartidas para las presentaciones Beamer en inglés (Metropolis)

% Conjuntos numéricos básicos
\newcommand{\R}{\mathbb{R}}
\newcommand{\N}{\mathbb{N}}
\newcommand{\Q}{\mathbb{Q}}
\newcommand{\Z}{\mathbb{Z}}
\newcommand{\set}[1]{\left\{ #1 \right\}}
\newcommand{\abs}[1]{\left|#1\right|}
\newcommand{\norm}[1]{\left\| #1 \right\|}

% Comandos estadísticos y probabilísticos del curso
\newcommand{\Var}{\operatorname{Var}}
\newcommand{\cov}{\operatorname{Cov}}
\newcommand{\corr}{\rho}
\newcommand{\comb}[2]{\begin{pmatrix} #1 \\ #2 \end{pmatrix}}
\newcommand{\s}{\sigma}
\newcommand{\particion}{\mathfrak{P}}
\newcommand{\card}[1]{\#\left( #1 \right)}
\newcommand{\seq}[3]{\set{#1}_{#2}^{#3}}
\newcommand{\E}{\mathbb{E}}
\newcommand{\Prob}{\mathbb{P}}

% Entornos pedagógicos y cajas en inglés académico natural (soportando tanto nombres en inglés como en español)
\newenvironment{discussion}{%
  \begin{alertblock}{Discussion Question}%
}{%
  \end{alertblock}%
}
\newenvironment{discusion}{%
  \begin{alertblock}{Discussion Question}%
}{%
  \end{alertblock}%
}

\newenvironment{reflection}{%
  \begin{alertblock}{Classroom Reflection}%
}{%
  \end{alertblock}%
}
\newenvironment{reflexion}{%
  \begin{alertblock}{Classroom Reflection}%
}{%
  \end{alertblock}%
}

\newenvironment{paradox}{%
  \begin{alertblock}{Exploratory Paradox}%
}{%
  \end{alertblock}%
}
\newenvironment{paradoja}{%
  \begin{alertblock}{Exploratory Paradox}%
}{%
  \end{alertblock}%
}

\newenvironment{motivation}{%
  \begin{exampleblock}{Motivation: Data Science in Practice}%
}{%
  \end{exampleblock}%
}
\newenvironment{motivacion}{%
  \begin{exampleblock}{Motivation: Data Science in Practice}%
}{%
  \end{exampleblock}%
}

\newenvironment{quickchallenge}{%
  \begin{block}{Quick Team Challenge}%
}{%
  \end{block}%
}
\newenvironment{retoclase}{%
  \begin{block}{Quick Team Challenge}%
}{%
  \end{block}%
}


\title[Categorical Variables]{Section 09.11: Categorical Variables and Dummy Variables}
\subtitle{Dummy Encoding, the $(n-1)$ Rule, and the Dummy Variable Trap}
\author[J. Castillo Colmenares]{Juliho Castillo Colmenares}
\institute{Tecnológico de Monterrey}
\date{\vspace{-1.5cm}}

\begin{document}

% Slide 1: Title
\begin{frame}[plain]
  \titlepage
\end{frame}

% Slide 2: Roadmap
\begin{frame}{Roadmap --- Chapter 09: Linear and Multiple Regressions}
  \scriptsize
  Where are we in the modular development of this chapter?
  \begin{itemize}\itemsep=0.02em
    \item \textbf{09.01} Correlation as the Premise of Regression
    \item \textbf{09.02} Introduction to Linear Regression
    \item \textbf{09.03} The Mathematics of Regression: OLS Derivation and $R^2$
    \item \textbf{09.04} Linear Regression on Simulated Data
    \item \textbf{09.05} Optimal Coefficients, $t$/$F$ Tests, and RSE
    \item \textbf{09.06} Implementation in Python with \texttt{statsmodels}
    \item \textbf{09.07} Multiple Regression, the Normal Equation, and Ridge/Lasso
    \item \textbf{09.08} Model Validation and $k$-fold Cross-Validation
    \item \textbf{09.09} Residual Diagnostics and Multicollinearity (VIF)
    \item \textbf{09.10} Regression with \texttt{scikit-learn}
    \item \textbf{09.11} Categorical Variables and Dummy Variables \emph{(Today)}
    \item \textbf{09.12} Nonlinear Transformations and Polynomial Regression
  \end{itemize}
  \pause
  \vspace{-0.05cm}
  \begin{block}{Session Objective}
    Incorporate categorical predictors (gender, city category) into a regression model via dummy variables; apply the $(n-1)$ rule to avoid the Dummy Variable Trap; and discover the exact mathematical equivalence between this technique and the one-way ANOVA from Chapter 08.
  \end{block}
\end{frame}

% Slide 3: Motivation
\begin{frame}{The Problem of Non-Numeric Predictors}
  \footnotesize
  \begin{motivacion}
    So far every predictor has been numeric. But many real datasets include categorical variables: gender, product category, geographic region. Assigning arbitrary codes (Female$=0$, Male$=1$, or worse: Tier1$=1$, Tier2$=2$, Tier3$=3$) would impose a \textbf{fictitious} order and numeric scale on categories that have none.
  \end{motivacion}

  \vspace{0.15cm}
  \pause
  The standard solution: represent each category via \textbf{binary indicator variables} (dummy variables).
\end{frame}

% Slide 4: Two-level encoding
\begin{frame}{Encoding a Two-Level Variable}
  \footnotesize
  \begin{block}{Definition}
    For a 2-level variable (e.g., Gender), we define a single indicator variable:
    $$D=\begin{cases}1 & \text{if the category of interest (e.g., Male)}\\0 & \text{if the reference category (Female)}\end{cases}$$
  \end{block}
  \pause

  \vspace{0.1cm}
  \begin{alertblock}{Model}
    $Y=\beta_0+\beta_1X+\beta_2D+\varepsilon$: the coefficient $\beta_2$ measures the average difference between categories, holding $X$ fixed.
  \end{alertblock}
\end{frame}

% Slide 5: The (n-1) rule
\begin{frame}{The $(n-1)$ Rule for Multi-Level Variables}
  \footnotesize
  \begin{block}{General Rule}
    For a categorical variable with $k$ levels, exactly $(k-1)$ dummy variables are created. One level remains the \textbf{base reference}, implicitly represented when all dummy variables are simultaneously zero.
  \end{block}
  \pause

  \vspace{0.1cm}
  \begin{alertblock}{Why Not $k$ Variables?}
    Including all $k$ dummy variables alongside the intercept produces the \textbf{Dummy Variable Trap}: perfect multicollinearity, since $\sum_{j=1}^kD_j=1$ for every observation.
  \end{alertblock}
\end{frame}

% Slide 6: Connection to ANOVA
\begin{frame}{A Deep Connection: Dummy Variables and ANOVA}
  \footnotesize
  \begin{block}{A Remarkable Discovery}
    One-way ANOVA (Section 08.01, $Y_{ij}=\mu+\tau_i+\epsilon_{ij}$) and regression with $(k-1)$ dummy variables are \textbf{mathematically identical}: $\beta_0=\mu_k$ (reference mean), $\beta_j=\mu_j-\mu_k$.
  \end{block}
  \pause

  \vspace{0.1cm}
  \begin{alertblock}{Same $F$-Test}
    The Partial $F$-Test on the dummy variable coefficients ($H_0:\beta_1=\cdots=\beta_{k-1}=0$) is \textbf{numerically identical} to the global $F$-Test of one-way ANOVA.
  \end{alertblock}
\end{frame}

% Slide 7: Python Lab (Block 1)
\begin{frame}[fragile]{Python Lab: Two-Level Variable}
  \scriptsize
  Manually encoding \texttt{Gender} as a dummy variable (\texttt{09.11\_categorical\_dummy\_variables.py}):
  {\lstset{basicstyle=\fontsize{3.4pt}{4.1pt}\selectfont\ttfamily,aboveskip=0.05em,belowskip=0.05em}
  \lstinputlisting[firstline=17, lastline=37]{../../code/09_regresiones/09.11_categorical_dummy_variables.py}}
\end{frame}

% Slide 8: Python Lab (Block 2)
\begin{frame}[fragile]{Python Lab: Three-Level Variable}
  \scriptsize
  Encoding \texttt{City Tier} (3 levels) with exactly 2 dummy variables:
  {\lstset{basicstyle=\fontsize{3.0pt}{3.7pt}\selectfont\ttfamily,aboveskip=0.05em,belowskip=0.05em}
  \lstinputlisting[firstline=39, lastline=61]{../../code/09_regresiones/09.11_categorical_dummy_variables.py}}
\end{frame}

% Slide 9: Python Lab (Block 3)
\begin{frame}[fragile]{Python Lab: the Dummy Variable Trap}
  \scriptsize
  Numeric verification: including all 3 dummy variables makes the design matrix singular:
  {\lstset{basicstyle=\fontsize{2.8pt}{3.4pt}\selectfont\ttfamily,aboveskip=0.05em,belowskip=0.05em}
  \lstinputlisting[firstline=64, lastline=86]{../../code/09_regresiones/09.11_categorical_dummy_variables.py}}
\end{frame}

% Slide 10: Terminal Output
\begin{frame}[fragile]{Python Lab: Terminal Output}
  \scriptsize
  Standard output produced when running the lab script:
  \vspace{0.02cm}
  \begin{block}{Standard Console Output}
    \fontsize{3.2pt}{3.9pt}\selectfont
    \begin{verbatim}
=== Block 1: Two-Level (Gender) ===
Male dummy coefficient: 231.41 (male spends $231.41 more, holding income fixed)

=== Block 2: Three-Level (City Tier), n-1 Rule ===
Tier1=390.72, Tier2=125.35 (baseline: Tier 3); 4 columns for 3 levels, no trap

=== Block 3: Dummy Variable Trap ===
Correct encoding (2 dummies):  condition number = 3.57
Trap encoding (all 3 dummies): condition number = 2.92e+15
Rank of trap matrix: 3 (columns: 4, rank-deficient: True)
    \end{verbatim}
  \end{block}
\end{frame}


% Slide 15: Closing
\begin{frame}{Section Closing: Categorical Variables}
  \begin{reflexion}
    Dummy variables extend regression to categorical predictors without imposing an artificial order, strictly following the $(k-1)$ rule to avoid the Dummy Variable Trap. Their exact equivalence with one-way ANOVA reveals that both techniques --- seemingly distinct in Chapters 08 and 09 --- are really the same statistical tool viewed from two angles.
  \end{reflexion}

  \vspace{0.2cm}
  \pause
  \begin{block}{Key Takeaways}
    \begin{itemize}\itemsep=0.08cm
      \item \textbf{$(k-1)$ rule:} $k$ levels require exactly $k-1$ dummy variables plus the intercept.
      \item \textbf{Dummy Variable Trap:} including all $k$ produces perfect collinearity (rank deficiency).
      \item \textbf{ANOVA $\equiv$ dummy regression:} same mathematics, same $F$-statistic, different notation.
    \end{itemize}
  \end{block}
\end{frame}

% Slide 16: Modular Perspective
\begin{frame}{Modular Perspective --- The Next Challenge}
  \scriptsize
  \begin{retoclase}
    We have handled \textbf{categorical} nonlinearity (qualitative variables). Section 09.12, the chapter's closing section, addresses \textbf{functional} nonlinearity: when the relationship between a numeric predictor and the response is not a straight line but a curve (quadratic, polynomial), showing that the same multiple regression engine from Section 09.07 solves this case too.
  \end{retoclase}

  \vspace{0.08cm}
  \pause
  \begin{alertblock}{Recommended Preparation}
    \begin{itemize}\itemsep=0.06cm
      \item Reflect on how you might include $X^2$ as a "new variable" within the already-familiar design matrix.
      \item Review the geometric interpretation of a parabola: vertex, concavity, and its meaning in a predictive model.
    \end{itemize}
  \end{alertblock}
\end{frame}

\end{document}
